\documentclass[12pt]{article}
\usepackage[margin=0.8in]{geometry}
\usepackage{amsmath}
\usepackage{xcolor}
\usepackage{tcolorbox}
\usepackage{enumitem}

\title{\textbf{Where Does $p \geq \frac{1}{2}$ Come From?}}
\author{}
\date{}

\begin{document}
\maketitle

\section*{The Setup}

Let me show you where the $p \geq \frac{1}{2}$ condition actually comes from by deriving the FOC.

\subsection*{The Problem}

You have initial wealth $w_0$. You choose how much $x$ to invest in a risky asset:

\begin{itemize}
\item With probability $p$: risky asset succeeds, you get back $Rx$ (where $R > 1$ is the return)
\item With probability $(1-p)$: risky asset fails, you get back $0$
\item The remaining $(w_0 - x)$ stays safe
\end{itemize}

Your final wealth is:
\begin{itemize}
\item \textbf{Good state (prob $p$):} $w_G = w_0 - x + Rx = w_0 + (R-1)x$
\item \textbf{Bad state (prob $1-p$):} $w_B = w_0 - x$
\end{itemize}

Your utility: CRRA with $u(w) = \frac{w^{1-\sigma}}{1-\sigma}$

\section*{Setting Up Expected Utility}

\begin{align*}
EU(x) &= p \cdot u(w_G) + (1-p) \cdot u(w_B) \\
&= p \cdot \frac{[w_0 + (R-1)x]^{1-\sigma}}{1-\sigma} + (1-p) \cdot \frac{[w_0 - x]^{1-\sigma}}{1-\sigma}
\end{align*}

\section*{Taking the First-Order Condition}

Take derivative with respect to $x$ and set equal to zero:

\begin{align*}
\frac{dEU}{dx} &= p \cdot \frac{[w_0 + (R-1)x]^{-\sigma}}{1-\sigma} \cdot (R-1) \\
&\quad + (1-p) \cdot \frac{[w_0 - x]^{-\sigma}}{1-\sigma} \cdot (-1) = 0
\end{align*}

Cancel out $\frac{1}{1-\sigma}$ (assuming $\sigma \neq 1$):

\begin{align*}
p(R-1)[w_0 + (R-1)x]^{-\sigma} &= (1-p)[w_0 - x]^{-\sigma}
\end{align*}

\section*{Rearranging}

\begin{align*}
\frac{p(R-1)}{1-p} &= \frac{[w_0 - x]^{-\sigma}}{[w_0 + (R-1)x]^{-\sigma}} \\
\frac{p(R-1)}{1-p} &= \left(\frac{w_0 + (R-1)x}{w_0 - x}\right)^{\sigma}
\end{align*}

Raise both sides to power $\frac{1}{\sigma}$:

\begin{align*}
\left[\frac{p(R-1)}{1-p}\right]^{1/\sigma} &= \frac{w_0 + (R-1)x}{w_0 - x}
\end{align*}

\newpage

\section*{Solving for $x$}

Let $K = \left[\frac{p(R-1)}{1-p}\right]^{1/\sigma}$ for simplicity.

\begin{align*}
K(w_0 - x) &= w_0 + (R-1)x \\
Kw_0 - Kx &= w_0 + (R-1)x \\
Kw_0 - w_0 &= Kx + (R-1)x \\
w_0(K - 1) &= x[K + (R-1)] \\
\end{align*}

\begin{tcolorbox}[colback=yellow!10!white,colframe=orange!75!black]
\textbf{The Solution:}
\[
\boxed{x^* = \frac{w_0(K - 1)}{K + (R-1)}}
\]

where $K = \left[\frac{p(R-1)}{1-p}\right]^{1/\sigma}$
\end{tcolorbox}

\section*{When is $x^* \geq 0$?}

For $x^* \geq 0$, we need:
\[
\frac{w_0(K - 1)}{K + (R-1)} \geq 0
\]

Since $w_0 > 0$ and $K + (R-1) > 0$ always (both positive), we need:
\[
K - 1 \geq 0 \quad \Rightarrow \quad K \geq 1
\]

That is:
\[
\left[\frac{p(R-1)}{1-p}\right]^{1/\sigma} \geq 1
\]

Raise both sides to power $\sigma$ (assuming $\sigma > 0$):
\[
\frac{p(R-1)}{1-p} \geq 1
\]

\begin{align*}
p(R-1) &\geq 1-p \\
p(R-1) + p &\geq 1 \\
pR &\geq 1 \\
\end{align*}

\begin{tcolorbox}[colback=green!5!white,colframe=green!75!black]
\textbf{General Condition:}
\[
\boxed{p \geq \frac{1}{R}}
\]

For $x^* \geq 0$, you need the probability of success to be at least $\frac{1}{R}$.
\end{tcolorbox}

\newpage

\section*{The Special Case: $R = 2$}

In your specific problem, if the risky asset \textbf{doubles your money} (pays back \$2 for every \$1 invested), then $R = 2$.

Substituting:
\[
p \geq \frac{1}{2}
\]

\textbf{That's where the $p \geq \frac{1}{2}$ comes from!}

\section*{Intuition}

\begin{tcolorbox}[colback=blue!5!white,colframe=blue!75!black]
\textbf{General Intuition:}

For the risky asset to be worth investing in at all, its \textbf{expected gross return} must be at least 1:

\[
pR + (1-p) \cdot 0 \geq 1
\]

This gives: $pR \geq 1$, or $p \geq \frac{1}{R}$

\vspace{0.3cm}

\textbf{In words:} "The probability of success times the payoff must be at least 1 for me to break even in expectation."
\end{tcolorbox}

\subsection*{Examples}

\textbf{Example 1: $R = 2$ (doubles your money)}
\begin{itemize}
\item Need $p \geq \frac{1}{2}$ (succeeds at least 50\% of the time)
\item Expected return: $2p \geq 1$
\end{itemize}

\textbf{Example 2: $R = 4$ (quadruples your money)}
\begin{itemize}
\item Need $p \geq \frac{1}{4}$ (succeeds at least 25\% of the time)
\item Even though success rate is lower, the bigger payoff ($R=4$) compensates
\end{itemize}

\textbf{Example 3: $R = 1.5$ (50\% return)}
\begin{itemize}
\item Need $p \geq \frac{1}{1.5} = \frac{2}{3} \approx 0.67$ (succeeds at least 67\% of the time)
\item Smaller payoff means you need higher success probability
\end{itemize}

\newpage

\section*{What if $p < \frac{1}{R}$?}

If $p < \frac{1}{R}$, then the FOC would give $x^* < 0$.

\textbf{What does negative $x$ mean?}
\begin{itemize}
\item Mathematically: you want to "invest negative amount" = short sell the risky asset
\item In reality: you can't do this (constraint $x \geq 0$)
\end{itemize}

So when $p < \frac{1}{R}$, the optimal choice given the constraint is:
\[
x^* = 0 \quad \text{(corner solution)}
\]

\textbf{Intuition:} The risky asset is such a bad bet (low probability of success relative to the payoff) that you shouldn't invest at all.

\section*{Summary}

\begin{tcolorbox}[colback=yellow!20!white,colframe=orange!75!black]
\textbf{The Complete Picture:}

\vspace{0.3cm}

\textbf{General condition:} $p \geq \frac{1}{R}$ for $x^* \geq 0$

\vspace{0.3cm}

\textbf{Your specific problem:} If $R = 2$, then $p \geq \frac{1}{2}$

\vspace{0.3cm}

\textbf{Decision rule:}
\begin{itemize}
\item If $p \geq \frac{1}{R}$: Use the FOC to get interior solution $x^* > 0$
\item If $p < \frac{1}{R}$: Corner solution, set $x^* = 0$ (don't invest)
\end{itemize}

\vspace{0.3cm}

\textbf{Why?} When $p < \frac{1}{R}$, the expected gross return $pR < 1$, meaning you expect to lose money on average. A risk-averse person definitely won't invest!
\end{tcolorbox}

\section*{Check Your Understanding}

\textbf{Q: Why do we need $p \geq \frac{1}{R}$ and not just $pR > 0$?}

A: Because we also keep $(w_0 - x)$ safe. The risky asset needs to beat the \textbf{outside option} of keeping money safe (which has return 1). So we need:
\[
\text{Expected return of risky asset} = pR \geq 1
\]

\vspace{0.3cm}

\textbf{Q: Does this mean if $p \geq \frac{1}{R}$, everyone invests the same amount?}

A: No! The exact amount $x^*$ depends on:
\begin{itemize}
\item $p$ (higher $p$ $\Rightarrow$ invest more)
\item $\sigma$ (higher risk aversion $\Rightarrow$ invest less)
\item $R$ (higher return $\Rightarrow$ invest more)
\item $w_0$ (more initial wealth $\Rightarrow$ invest more in absolute terms)
\end{itemize}

\end{document}
